\documentclass[10pt,a4paper]{amsart}
\usepackage[margin=19mm]{geometry}
\usepackage{amsmath,amssymb,mathtools}
\usepackage[scr=rsfs,cal=euler]{mathalfa}
\usepackage{xfrac}
\usepackage[unicode,hidelinks,bookmarks=false]{hyperref}
\newcommand{\ie}{i.e.\ }
\newcommand{\cf}{cf.\ }
\newcommand{\resp}{respectively\ }
\newcommand{\Subsection}{\subsection}
\providecommand{\nobreakdash}{\nobreak\hbox{-}}
\setlength{\emergencystretch}{2em}
\multlinegap0pt
\usepackage{amssymb}
\usepackage{enumerate}
\usepackage{xcolor}
\usepackage{algorithm}\usepackage{algpseudocode}
\usepackage{bm}
\usepackage{mathtools}
\newtheorem{lemma}{Lemma}
\title{Stable Hermite transforms via the Golub--Welsch algorithm}
\author{Marcus Webb, Georg Maierhofer}
\date{Source: arXiv:2604.02041v2 (CC BY 4.0)}
\begin{document}
\maketitle

\noindent\emph{Source and licence.} Stable Hermite transforms via the Golub--Welsch algorithm. Version arXiv:2604.02041v2 (CC BY 4.0), distributed by arXiv under the Creative Commons Attribution 4.0 International licence, http://creativecommons.org/licenses/by/4.0/, as recorded in the arXiv licence metadata for this exact version and shown on the abstract page https://arxiv.org/abs/2604.02041v2. Full licence notice: \texttt{licenses/arxiv-2604.02041-CC-BY-4.0.txt}.\par\smallskip

\noindent\textit{Editorial scope.} Counted unit: the article's lemma lem:orth (source0.tex lines 265--274). The article's complete original proof of it is delivered with it (source0.tex lines 276--278, one \texttt{proof} environment of 3 lines). No mathematics is written, repaired or re-derived: the statement and the proof are the article's own lines, in the article's own notation, and no retained line is substituted or re-flowed.  Retained uncounted as the source-local context the proof needs: lines 210--212; lines 220--222.  Nothing was omitted from any retained span, so the package carries no omission record.  No retained line contains a citation, so the wrapper carries no bibliography and no bibliography entry.  No retained line contains a figure environment, an image-inclusion command or drawing code, so no image is delivered and the package declares no asset dependency.  Editorial formatting only, in addition: an AMS \texttt{amsart} preamble that restates the article's own declarations and macros used by the retained lines, namely the environments \texttt{lemma} on separate counters, together with the standard packages those lines need.  The retained environments print in this document as Lemma 1 in order of appearance, whereas the article prints the same items with the numbering of its own full document.  The complete native source of the article as distributed in the arXiv e-print for this exact version is bundled unchanged as \texttt{sources/197702/source0.tex}.
\begin{equation}\label{eqn:Tinv}
    T^{-1} = T^T W,
\end{equation}

\begin{equation}\label{eqn:ghweights}
    w_k = \frac{\mathrm{e}^{-x_k^2}}{N \psi_{N{-}1}(x_k)^2}.
\end{equation}

\begin{lemma}\label{lem:orth}
The Hermite transform matrix $T \in \mathbb{R}^{N\times N}$, with the nodes $\mathbf{x}$ taken to be the Gauss--Hermite quadrature nodes, satisfies
\begin{equation}
    T = DQ^T, \qquad T^{-1} = Q D^{-1},
\end{equation}
where $Q \in \mathbb{R}^{N\times N}$ is orthogonal and $D \in \mathbb{R}^{N\times N}$ is diagonal, with
\begin{align*}
    {Q}_{kj} = \frac{\psi_k(x_j)}{\sqrt{N} \left|\psi_{N{-}1}(x_j)\right|}, \qquad {D}_{jj} = \sqrt{N}\left|\psi_{N{-}1}(x_j)\right|.
\end{align*}
\end{lemma}

\begin{proof}
Clearly $T = DQ^T$ by the definition of $D$ and $Q$. What remains to show is that $Q$ is orthogonal. Consider equation \eqref{eqn:Tinv}. It says that ${T}^T {W} {T} = {I}$, where $W$ is a diagonal matrix with elements $W_{kk} = N^{-1}\psi_{N{-}1}(x_k)^{-2}$ by equation \eqref{eqn:ghweights}. This implies that $W^{1/2}T = Q$ is orthogonal.
\end{proof}

\end{document}
